% =====================================================================
% appendices/appendix_a_implementation.tex
% Appendix A — Implementation Details
% Project: Vision-Based Running Technique Classification from In-the-Wild
%          Internet Videos using Markerless Pose Estimation and Deep Learning
%
% Required preamble packages in main.tex:
%   \usepackage{booktabs}
%   \usepackage{array}
%   \usepackage{tabularx}
%   \usepackage{graphicx}
%
% This appendix documents implementation only. Configuration values are
% reported only when available from the project notebooks, exported reports,
% or experiment inventory. Missing training hyper-parameters are explicitly
% marked as not exported rather than inferred.
% =====================================================================

\chapter{Implementation Details}
\label{app:implementation}

This appendix documents the implementation behind the experiments of
Chapters~\ref{ch:experimental-setup}--\ref{ch:results}: the notebook
organisation, the output directory layout, the files each stage produces, and
the configuration settings. It is provided for reproducibility; no additional
result claims are introduced here. Configuration values are reported only when
they are available from the project notebooks or exported reports; values that
were not exported are explicitly marked as unavailable rather than inferred.

% ---------------------------------------------------------------------
\section{Notebook Mapping}
\label{app:notebook-mapping}

The pipeline is implemented as a sequence of self-contained notebooks.
Table~\ref{tab:notebook-mapping} maps each notebook to the pipeline stage it
realises and the principal role it plays in the thesis. Notebooks numbered with
a letter suffix (03b--03g) are extended or supplementary analyses; the others
are core pipeline stages.

\begin{table}[htbp]
\centering
\caption{Mapping of project notebooks to pipeline stages and thesis roles.}
\label{tab:notebook-mapping}
\renewcommand{\arraystretch}{1.18}
\scriptsize
\begin{tabularx}{\textwidth}{
>{\raggedright\arraybackslash}p{0.08\textwidth}
>{\raggedright\arraybackslash}p{0.34\textwidth}
>{\raggedright\arraybackslash}X
>{\raggedright\arraybackslash}p{0.16\textwidth}}
\toprule
ID & Notebook & Purpose & Role \\
\midrule
00  & \path{00_dataset_split_and_eda} & Dataset construction, video-level splitting, and exploratory data analysis. & Core \\
01  & \path{01_handmade_8frame_upper_bound} & Handmade eight-phase cycle pipeline used as the controlled upper-bound reference. & Core \\
02  & \path{02_autocorr_cycle_detection_baseline} & Automatic autocorrelation-based landmark baseline using motion signals and detected cycles. & Core \\
03  & \path{03_handmade_vs_autocorr_comparison} & Comparison between manually controlled cycle segmentation and automatic cycle detection. & Core \\
03b & \path{03b_autocorr_image_based_pipeline} & Image-based autocorr pipeline using runner crops and a MobileNetV2--LSTM classifier. & Supplementary \\
03c & \path{03c_autocorr_version_comparison} & Autocorr version history and method-evolution summary. & Supplementary \\
03d & \path{03d_representation_comparison_landmark_vs_image} & Landmark-only versus image-based representation comparison. & Supplementary \\
03e & \path{03e_cycle_level_vs_video_level_evaluation} & Cycle-level and video-level evaluation with aggregation. & Supplementary \\
03f & \path{03f_pose_cycle_quality_analysis} & Pose-quality and cycle-quality analysis for interpreting failures. & Supplementary \\
03g & \path{03g_threshold_05_vs_tuned_report} & Comparison between default threshold 0.5 and validation-tuned threshold. & Supplementary \\
03h & \path{03h_autocorr_cv_robustness_check} & Fair handmade--autocorr gap under 5-fold video-level CV and autocorr MLP permutation test. & Core support \\
04  & \path{04_model_ablation} & Architecture ablation over MLP, 1D-CNN, BiLSTM, and CNN-BiLSTM. & Core \\
05  & \path{05_feature_ablation} & Input feature-set comparison. & Core \\
06  & \path{06_signal_ablation_cycle_detection} & Comparison of motion signals for cycle detection. & Core \\
07  & \path{07_frames_per_cycle_ablation} & Comparison of 8, 16, 24, and 32 sampled frames per cycle. & Core \\
08  & \path{08_boundary_detector_pipeline} & Learned boundary-detection pipeline. & Supplementary \\
09  & \path{09_raw_sequence_bilstm_gru} & Raw-sequence BiLSTM/GRU pipeline without fixed eight-phase sampling. & Supplementary \\
10  & \path{10_threshold_tuning_and_calibration} & Threshold sweep and calibration analysis. & Supplementary \\
11  & \path{11_error_analysis} & Error analysis used to support the discussion of failure modes. & Core support \\
12  & \path{12_final_report_tables_and_figures} & Core report tables and figures. & Core \\
12b & \path{12b_extended_final_report_tables} & Extended report tables and final summary figures. & Supplementary \\
\bottomrule
\end{tabularx}
\end{table}

% ---------------------------------------------------------------------
\section{Directory Structure}
\label{app:directory}

All notebooks write to a fixed output directory under a configurable
\texttt{DRIVE\_BASE} root, so that downstream aggregation notebooks can locate
artefacts regardless of where the project is mounted. The principal layout is
shown below.

\lstset{
  breaklines=true,      % Tự động ngắt dòng khi chạm lề
  basicstyle=\ttfamily\small, % Phông chữ máy tính nhỏ gọn gọn gàng
  columns=flexible,
  keepspaces=true
}

\begin{lstlisting}
DRIVE_BASE/
  outputs/
    dataset/          # split lists, split_summary, metadata
    eda/              # class/duration/fps/frame distributions, split charts
    handmade/         # handmade results, predictions, reports, curves, *.pth
    autocorr/         # autocorr results, cycle stats, fair-gap CV/permutation files
    autocorr_image/   # image results, cycle stats, crops, reports, *.pth
    ablation/         # model / feature / signal / frames ablation CSVs
    reports/          # comparison CSVs/MD, extended summaries, figures
    figures/          # copied thesis-ready figures
\end{lstlisting}

% ---------------------------------------------------------------------
\section{Output Files}
\label{app:output-files}

Table~\ref{tab:output-files} lists the principal machine-readable outputs and
the notebook that produces each. These files are the inputs to the controlled
comparisons and ablations of Chapter~\ref{ch:results}.

\begin{table}[htbp]
\centering
\caption{Principal output files and producing notebooks.}
\label{tab:output-files}
\renewcommand{\arraystretch}{1.18}
\scriptsize
\begin{tabularx}{\textwidth}{
>{\raggedright\arraybackslash}p{0.38\textwidth}
>{\raggedright\arraybackslash}p{0.16\textwidth}
>{\raggedright\arraybackslash}X}
\toprule
File & Producing notebook & Contents \\
\midrule
\texttt{dataset/split\_summary.csv} & 00 & Per-split video and cycle counts. \\
\texttt{handmade/handmade\_results.csv} & 01 & Handmade upper-bound metrics for the landmark models. \\
\texttt{handmade/handmade\_predictions.csv} & 01 & Per-sample predictions for the handmade upper-bound pipeline. \\
\texttt{handmade/handmade\_classification\_report.txt} & 01 & Classification report for the handmade pipeline. \\
\texttt{autocorr/autocorr\_results.csv} & 02 & Autocorr landmark baseline metrics. \\
\texttt{autocorr/autocorr\_predictions.csv} & 02 & Per-sample predictions for the autocorr landmark pipeline. \\
\texttt{autocorr/autocorr\_cycle\_stats.csv} & 02 & Per-cycle detection statistics and quality information. \\
\texttt{autocorr/autocorr\_failed\_videos.csv} & 02 & Failed or rejected videos in the autocorr pipeline. \\
\texttt{autocorr/fair\_gap\_and\_mlp\_perm.csv} & 03h & Fair-gap CV summary and autocorr MLP permutation-test values. \\
\texttt{autocorr/fair\_gap\_and\_mlp\_perm.json} & 03h & Machine-readable fair-gap and permutation-test record. \\
\texttt{reports/pipeline\_comparison.csv} & 03 & Handmade versus autocorr comparison. \\
\texttt{autocorr\_image/image\_results.csv} & 03b & Image-based autocorr pipeline metrics. \\
\texttt{autocorr\_image/image\_predictions.csv} & 03b & Per-sample predictions for the image-based pipeline. \\
\texttt{autocorr\_image/image\_cycle\_stats.csv} & 03b & Cycle and crop statistics for the image pipeline. \\
\texttt{reports/autocorr\_version\_comparison.csv} & 03c & Autocorr version history and method-evolution records. \\
\texttt{reports/representation\_comparison.csv} & 03d & Landmark versus image representation comparison. \\
\texttt{reports/cycle\_level\_metrics.csv} & 03e & Cycle-level evaluation metrics. \\
\texttt{reports/video\_level\_metrics.csv} & 03e & Video-level aggregated evaluation metrics. \\
\texttt{reports/pose\_cycle\_quality\_metrics.csv} & 03f & Quality-group metrics for pose and cycle quality analysis. \\
\texttt{reports/threshold\_05\_vs\_tuned.csv} & 03g & Default-threshold versus tuned-threshold comparison. \\
\texttt{ablation/model\_ablation.csv} & 04 & Architecture ablation results. \\
\texttt{ablation/feature\_ablation.csv} & 05 & Feature-set ablation results. \\
\texttt{ablation/signal\_ablation.csv} & 06 & Motion-signal ablation results. \\
\texttt{ablation/frames\_ablation.csv} & 07 & Frames-per-cycle ablation results. \\
\texttt{bd\_results.csv} & 08 & Boundary-detector pipeline results, when exported. \\
\texttt{raw\_sequence\_results.csv} & 09 & Raw-sequence BiLSTM/GRU results, when exported. \\
\texttt{reports/main\_pipeline\_comparison.csv} & 12b & Consolidated main pipeline summary. \\
\texttt{reports/extended\_final\_results\_summary.csv} & 12b & Extended final summary including supplementary pipelines. \\
\bottomrule
\end{tabularx}
\end{table}

% ---------------------------------------------------------------------

\section{Hyperparameter Settings}
\label{app:hyperparams}

Table~\ref{tab:hyperparams} records implementation settings that could be
verified from the project notebooks, exported reports, and the hyper-parameter
evidence audit. The table intentionally separates settings that were verified
from settings that were not exported. This prevents the thesis from inventing
training details while still documenting the values that are available in the
experiment evidence.

\begin{table}[htbp]
\centering
\caption{Verified implementation settings and model-configuration evidence.}
\label{tab:hyperparams}
\renewcommand{\arraystretch}{1.15}
\scriptsize
\begin{tabularx}{\textwidth}{
>{\raggedright\arraybackslash}p{0.31\textwidth}
>{\raggedright\arraybackslash}X}
\toprule
Setting & Verified value / status \\
\midrule
Pose estimator & MediaPipe BlazePose, producing 33 pose landmarks with landmark visibility. \\
Selected landmarks & 13 landmarks: 0, 11, 12, 23, 24, 25, 26, 27, 28, 29, 30, 31, and 32. \\
Cycle representation & Eight sampled phases per detected running cycle. \\
Velocity augmentation & Enabled for the landmark representation where selected landmark coordinates are augmented with temporal velocity features. \\
Main automatic cycle signal & \texttt{contact\_diff}. \\
Autocorr cycle logic & Foot-strike-to-foot-strike cycle detection with two-stride correction. \\
Autocorr visibility threshold & 0.45. \\
Autocorr cycle duration bounds & $[0.50, 1.15]$ seconds for the controlled landmark baseline; the historical image-development notebook used a wider automatic-detection range before stricter cycle filtering. \\
Autocorr fixed margin & Disabled; \texttt{USE\_MARGIN = False}. \\
Image crop size & $224 \times 224$ RGB runner-centred crop. \\
Image backbone & MobileNetV2 visual feature extractor with an LSTM temporal head. \\
\midrule
Landmark model training schedule & Verified from the audited landmark notebooks: learning rate $10^{-3}$, batch size 16, 80 epochs, weight decay $10^{-4}$, random seed 42. Optimiser and loss function were not consistently exported in the final evidence. \\
MLP parameter count & 193{,}025 parameters. \\
1D-CNN parameter count & 40{,}257 parameters. \\
BiLSTM parameter count & 173{,}441 parameters. \\
CNN-BiLSTM parameter count & 82{,}113 parameters. \\
\midrule
Autocorr image representation schedule & Verified from the audited \texttt{03b\_autocorr\_image\_based\_pipeline.ipynb}: learning rate $10^{-3}$, batch size 8, 40 epochs, weight decay $10^{-4}$, dropout 0.4, two recurrent layers, random seed 42. \\
Historical \texttt{autocorr\_fixed} image-development model & Verified from the historical 76.5\% notebook: 8 RGB frames per cycle, batch size 8, MobileNetV2 ImageNet backbone, LSTM hidden size 128, two LSTM layers, dropout 0.4, LayerNorm + two-layer classification head. \\
Historical \texttt{autocorr\_fixed} training schedule & AdamW with OneCycleLR in two phases: 15 frozen-backbone epochs with lr $3\times10^{-4}$ and weight decay $10^{-3}$, followed by 10 partial-fine-tuning epochs with lr $5\times10^{-5}$ and weight decay $10^{-4}$. The last five MobileNetV2 feature blocks were unfrozen in the fine-tuning phase. \\
Historical \texttt{autocorr\_fixed} loss and regularisation & \texttt{BCEWithLogitsLoss} with class \texttt{pos\_weight}, gradient clipping at 1.0, validation accuracy used for checkpoint selection in the historical development notebook. \\
MobileNetV2--LSTM parameter count & \textit{Not exported in provided reports}. \\
\midrule
Boundary-detector evidence & Dropout values 0.3/0.4 and two recurrent layers were found in the audit, but the full training schedule was not consistently exported; the boundary detector is therefore treated as supplementary. \\
GRU / raw-sequence parameter count & \textit{Not exported in provided reports}. \\
Raw-sequence training schedule & \textit{Not exported in provided reports}. \\
Early stopping & \textit{Not exported in provided reports}. \\
\bottomrule
\end{tabularx}

\vspace{0.5em}
\footnotesize
\noindent\textit{Note.} Parameter counts for MLP, 1D-CNN, BiLSTM, and
CNN-BiLSTM are taken from the exported model-ablation report. Landmark and image
training schedules are reported only when they were explicitly recovered from
the audited notebooks or the hyper-parameter evidence audit. The historical
\texttt{autocorr\_fixed} settings describe the development image pipeline that
produced the 76.5\% validation-accuracy milestone; they are documented as method
evidence and should not be confused with a controlled benchmark result unless
that pipeline is rerun under the shared train/validation/test protocol.
\end{table}

% ---------------------------------------------------------------------
\section{Reproducibility Notes}
\label{app:reproducibility-notes}

The experiments use video-level splitting to prevent cycles extracted from the
same source video from appearing in both training and evaluation partitions.
The handmade pipeline should be interpreted as an upper-bound reference because
the cycle boundaries are manually controlled. The autocorr landmark pipeline is
the main automatic baseline because it uses automatic cycle detection and the
fixed split used for the final thesis results. The image-based pipeline is an
extended representation comparison and should not be conflated with historical
image-based development notebooks unless rerun under the same fixed-split
protocol.

% End of Appendix A
